% =====================================================================
%  04-cli.tex - Interfaz de linea de comandos
% =====================================================================

\section{La interfaz de línea de comandos}\label{sec:impl-cli}

El primer adaptador driving entregado por el prototipo es la interfaz de línea de comandos \texttt{despacho-cli}, construida con \texttt{clap}~4 en modo derive (\rutarepo{crates/bin/src/cli.rs}). El diseño anticipó un adaptador de línea de comandos para tareas administrativas junto a los adaptadores REST y WebSocket (\cref{sec:vista-componentes}); el orden de construcción privilegió la CLI porque permitió ejercitar primero los servicios criptográficos de extremo a extremo y después reutilizarlos desde HTTP. La bandera global \texttt{-{}-json} conmuta las operaciones administrativas a un objeto JSON legible por máquina. La \cref{tab:impl-cli-catalogo} presenta el catálogo completo de subcomandos con sus banderas reales, y la \cref{tab:impl-cli-casos} el cableado de cada grupo hacia los casos de uso de \texttt{application} y los adaptadores concretos que la raíz de composición inyecta.

\begingroup
\footnotesize
\begin{longtable}{@{}L{7.2cm}L{\dimexpr\linewidth-7.2cm-2\tabcolsep\relax}@{}}
  \caption{Catálogo de subcomandos de \texttt{despacho-cli}. La bandera global \texttt{-{}-json} aplica a todos.}
  \label{tab:impl-cli-catalogo} \\
    \toprule
    \textbf{Comando y banderas} & \textbf{Efecto} \\
    \midrule
    \endfirsthead
    \toprule
    \textbf{Comando y banderas} & \textbf{Efecto} \\
    \midrule
    \endhead

    \texttt{serve -{}-signer-cert <CERT> -{}-signer-key <KEY> [-{}-bind <ADDR>] [-{}-data-dir <DIR>]} & Inicia la API por expediente con un rol PostgreSQL restringido, sesiones Redis y TSA local. \\
    \midrule
    \texttt{database migrate -{}-runtime-role <ROLE>} & Aplica esquema y permisos con una conexión administrativa; el rol debe existir. \\
    \texttt{database import [-{}-data-dir <DIR>] -{}-mapping <PATH> [-{}-apply]} & Inspecciona origen y destino; solo \texttt{-{}-apply} confirma la importación y el corte. \\
    \midrule
    \texttt{crypto hash <FILE>} & Resumen SHA-256 del archivo. \\
    \midrule
    \texttt{vault encrypt <FILE> -{}-doc-id <UUID> [-{}-version <N>]} & Cifra con envelope encryption; escribe \texttt{<FILE>.enc}. \\
    \texttt{vault decrypt <PACKAGE> -{}-doc-id <UUID> [-{}-version <N>] [-{}-out <PATH>]} & Descifra un archivo \texttt{DVLT1}; sin \texttt{-{}-out} escribe el claro a la salida estándar. \\
    \texttt{vault rotate-kek -{}-file <PATH>} & Re-envuelve la DEK bajo la KEK de \texttt{NEW\_KEK\_BASE64}. \\
    \midrule
    \texttt{pki [-{}-scripts-dir <DIR>] init-ca} & Crea la CA raíz interna (idempotente). \\
    \texttt{pki issue -{}-cn <CN> [-{}-purpose <PURPOSE>]} & Emite un certificado ordinario o para declaración personal interna. \\
    \texttt{pki issue-request -{}-csr <CSR> -{}-name <NAME>} & Emite un certificado de firma desde un CSR externo sin generar la clave privada del titular. \\
    \texttt{pki revoke -{}-serial <HEX>} & Revoca por número de serie. \\
    \texttt{pki gen-crl} & Genera la CRL en PEM. \\
    \texttt{pki show <CERT>} & Inspecciona sujeto, emisor, serie y vigencia. \\
    \texttt{credential-trust publish -{}-root-cert <PATH> -{}-crl <PATH> -{}-expected-revision <N>} & Publica raíz y CRL con revisión esperada y conexión administrativa. \\
    \midrule
    \texttt{sign <FILE> -{}-cert <PATH> -{}-key <PATH>} & Firma el digest del archivo; escribe \texttt{<FILE>.sig}. \\
    \texttt{timestamp <INPUT> [-{}-mock] [-{}-out <PATH>] [-{}-pki-dir <DIR>]} & Solicita un sello RFC 3161; escribe \texttt{<INPUT>.tsr}. Con \texttt{-{}-mock} usa la TSA local; sin ella, el proveedor remoto. \\
    \texttt{verify <FILE> -{}-sig <PATH> -{}-cert <PATH> -{}-ca <PATH> [-{}-tsr] [-{}-crl] [-{}-tsa-cert] [-{}-at <RFC3339>]} & Verificación integral de cuatro componentes; \texttt{-{}-at} evalúa el certificado a un instante dado. \\
    \texttt{package export <FILE> -{}-sig -{}-tsr -{}-cert -{}-ca -{}-crl [-{}-tsa-chain] -{}-out <PATH>} & Exporta el paquete de evidencia, previa verificación de los artefactos. \\
    \midrule
    \texttt{auth calibrate} & Mide el costo real de Argon2id contra la banda de 500--1\,000 ms. \\
    \texttt{auth hash-password} & Hash PHC de una contraseña leída de la entrada estándar. \\
    \texttt{auth verify-password -{}-hash <PHC>} & Verifica una contraseña contra un PHC almacenado. \\
    \texttt{auth totp enroll -{}-user <EMAIL> -{}-secret-out <PATH>} & Enrola el segundo factor y emite los códigos de recuperación. \\
    \texttt{auth totp verify -{}-code <CODE> -{}-secret-file <PATH>} & Verifica un código TOTP. \\
    \midrule
    \texttt{audit append -{}-action <STRING> -{}-resource <STRING>} & Anexa un evento a la bitácora offline, si el origen no está bloqueado por migración. \\
    \texttt{audit verify-chain} & Verifica la cadena del archivo offline desde el génesis. \\
    \texttt{audit show} & Lista los eventos del archivo offline en orden de inserción. \\
    \bottomrule
\end{longtable}
\endgroup

\begin{table}[htbp]
  \centering
  \caption{Correspondencia entre comandos, orquestación y adaptadores compuestos por el binario.}
  \label{tab:impl-cli-casos}
  \begin{tabularx}{\linewidth}{@{}L{3.4cm} L{5.2cm} >{\raggedright\arraybackslash}X@{}}
    \toprule
    \textbf{Comando} & \textbf{Caso(s) de uso} & \textbf{Adaptadores inyectados} \\
    \midrule
    \texttt{serve} & \texttt{Identity\allowbreak Service}, \texttt{CaseService}, \texttt{CaseDocument\allowbreak Service} & PostgreSQL transaccional, Redis, Argon2id, TOTP, cifrado, firma, TSA local y ZIP \\
    \texttt{database *} & Preparación administrativa y validación de registros históricos & \texttt{initialize\_database}, \texttt{Legacy\allowbreak Import} y puertos de validación criptográfica \\
    \texttt{crypto hash} & \texttt{Hash\allowbreak Document} & \texttt{Ring\allowbreak Sha256\allowbreak Hasher} \\
    \texttt{vault *} & \texttt{Encrypt\allowbreak Document}, \texttt{Decrypt\allowbreak Document}, \texttt{RotateKek} & \texttt{Ring\allowbreak Aes\allowbreak Gcm\allowbreak Cipher}, \texttt{Envelope\allowbreak Key\allowbreak Manager} \\
    \texttt{pki *} & \texttt{Initialize\allowbreak Ca}, \texttt{Issue\allowbreak Certificate}, \texttt{Revoke\allowbreak Certificate}, \texttt{GenerateCrl}, \texttt{Inspect\allowbreak Certificate} & \texttt{Openssl\allowbreak Ca\allowbreak Adapter}, \texttt{X509\allowbreak Chain\allowbreak Validator} \\
    \texttt{credential-trust} & Publicación auditada de confianza interna & PostgreSQL administrativo y verificador de declaraciones \\
    \texttt{sign} & \texttt{Sign\allowbreak Document} + \texttt{Verify\allowbreak Signature} & \texttt{Rsa\allowbreak Pkcs1\allowbreak Signer}, \texttt{Rsa\allowbreak Pkcs1\allowbreak Verifier} \\
    \texttt{timestamp} & \texttt{Timestamp\allowbreak Document} & \texttt{Local\allowbreak Openssl\allowbreak Tsa} o \texttt{Cincel\allowbreak Tsa\allowbreak Adapter} \\
    \texttt{verify} & \texttt{Verify\allowbreak Document} & \texttt{Ring\allowbreak Sha256\allowbreak Hasher}, \texttt{Rsa\allowbreak Pkcs1\allowbreak Verifier}, \texttt{X509\allowbreak Chain\allowbreak Validator}, \texttt{Rfc3161\allowbreak Verifier} \\
    \texttt{package export} & \texttt{Verify\allowbreak Document} + \texttt{Export\allowbreak Evidence\allowbreak Package} & los cuatro anteriores más \texttt{Stored\allowbreak Zip\allowbreak Writer} \\
    \texttt{auth calibrate}, \texttt{hash-password}, \texttt{verify-password} & \texttt{CalibratePassword\allowbreak Hashing}, \texttt{Hash\allowbreak Password}, \texttt{Verify\allowbreak Password} & \texttt{Argon2id\allowbreak Hasher} \\
    \texttt{auth totp *} & \texttt{EnrollTotp}, \texttt{VerifyTotp} & \texttt{Totp\allowbreak Rs\allowbreak Provider}, \texttt{Random\allowbreak Recovery\allowbreak Code\allowbreak Generator}, \texttt{Argon2id\allowbreak Hasher} \\
    \texttt{audit *} & \texttt{Append\allowbreak Audit\allowbreak Event}, \texttt{Verify\allowbreak Audit\allowbreak Chain}, \texttt{Show\allowbreak Audit\allowbreak Log} & \texttt{File\allowbreak Audit\allowbreak Log}, \texttt{Ring\allowbreak Sha256\allowbreak Hasher}, \texttt{SystemClock} \\
    \bottomrule
  \end{tabularx}
\end{table}

Con \texttt{-{}-json}, las operaciones finitas emiten resultados estructurados en la salida estándar ---por ejemplo, \texttt{crypto hash} emite \texttt{\{algorithm, digest, file\}} y \texttt{verify} emite el reporte completo con los cuatro componentes y el veredicto---, mientras los diagnósticos van a la salida de error. Se exceptúan el proceso persistente \texttt{serve} y \texttt{vault decrypt} sin archivo de salida, que conserva el texto claro crudo. Los códigos de salida están pensados para que los scripts bifurquen directamente: el proceso termina con código distinto de cero cuando \texttt{verify} produce un veredicto no válido, cuando \texttt{auth verify-password} reporta una contraseña que no coincide, cuando \texttt{auth totp verify} rechaza el código, cuando \texttt{audit verify-chain} encuentra la cadena rota ---imprimiendo antes el JSON de fallo con el índice del primer eslabón roto--- y cuando \texttt{package export} rehúsa exportar. Los subcomandos de lectura de la bitácora (\texttt{verify-chain}, \texttt{show}) rehúsan operar sobre un archivo inexistente: reportar una ruta equivocada o un log borrado como una cadena vacía válida frustraría el propósito de la verificación.

La raíz de composición (\rutarepo{crates/bin/src/main.rs}) es el único lugar del sistema donde los adaptadores concretos se instancian e inyectan en los casos de uso, y el único crate que usa \texttt{anyhow}, tal como estableció el patrón de inyección de dependencias por composición del diseño (\cref{sec:inyeccion-dependencias}). Ahí ---y no en los adaptadores, que quedan libres de acceso al entorno--- se leen los secretos: la KEK de \texttt{KEK\_BASE64} a un buffer \texttt{Zeroizing}, la clave privada del firmante desde su archivo PEM, la credencial del proveedor de sellado. La raíz concentra también salvaguardas operativas que no pertenecen a ningún servicio individual: \texttt{sign} verifica la firma recién producida contra el certificado \emph{antes} de escribir \texttt{<FILE>.sig}, de modo que una clave que no corresponde al certificado falla ahí y no produce una firma inverificable; \texttt{timestamp} relee el token recibido y comprueba que cubra el digest calculado antes de escribir el archivo \texttt{.tsr}; \texttt{vault rotate-kek} escribe a un archivo temporal y renombra encima, de modo que un fallo a la mitad deja el original intacto; y \texttt{package export} somete los artefactos al mismo caso de uso \texttt{Verify\allowbreak Document} de la verificación integral ---CRL y token incluidos--- y rehúsa exportar si el veredicto no es válido, nombrando cada componente fallido, porque un paquete que fallaría sus propias instrucciones de verificación no debe existir.

La integración de todos los servicios se demuestra con el guion \rutarepo{scripts/demo.sh}, que compila el binario y recorre el ciclo completo de evidencia en un directorio temporal: creación de la autoridad interna, del firmante y de la autoridad de sellado; hash del documento de muestra; cifrado en reposo con corrupción deliberada de un byte y observación del rechazo; firma y obtención del sello de tiempo; verificación integral con los cuatro componentes válidos; alteración de una copia del documento con el consecuente fallo del componente de firma; exportación del paquete de evidencia y su verificación con solo \texttt{openssl} y \texttt{unzip}; revocación del certificado y su detección; calibración de autenticación y códigos de un solo uso; y bitácora encadenada con anexado, verificación, manipulación y nueva verificación. La transcripción íntegra de una ejecución real, capturada el 12 de julio de 2026 sobre OpenSSL 3.5.7, está versionada en \rutarepo{docs/demo-transcript.md}. Una segunda demostración, \rutarepo{scripts/api-demo.sh}, recorre el mismo núcleo a través de HTTP y se describe en la sección siguiente. El entorno, los comandos y el resumen de ambos recorridos quedaron asentados en \rutarepo{docs/verification-report.md}.
